%------------------------------------------------引用的宏包和相应的定义-----------------------------------------
\RequirePackage{silence}
\WarningFilter{fontspec}{Font}
\ActivateWarningFilters

\usepackage{xeCJK}          % 中文支持宏包
\usepackage{newunicodechar}
\usepackage[export]{adjustbox}      

\iffalse
\setCJKmainfont[BoldFont={SimHei},ItalicFont={[simkai.ttf]}]{SimSun}
\setCJKsansfont{SimHei}
\setCJKmonofont{[simfang.ttf]}

\setCJKfamilyfont{zhsong}{SimSun}
\setCJKfamilyfont{zhhei}{SimHei}
\setCJKfamilyfont{zhkai}{[simkai.ttf]}
\setCJKfamilyfont{zhfs}{[simfang.ttf]}
\setCJKfamilyfont{zhli}{LiSu}
%\setCJKfamilyfont{zhyou}{YouYuan}


\newcommand*{\songti}{\CJKfamily{zhsong}}
\newcommand*{\heiti}{\CJKfamily{zhhei}}
\newcommand*{\kaishu}{\CJKfamily{zhkai}}
\newcommand*{\fangsong}{\CJKfamily{zhfs}}
\newcommand*{\lishu}{\CJKfamily{zhli}}
%\newcommand*{\youyuan}{\CJKfamily{zhyou}}
\fi

\usepackage[paperwidth=170mm,paperheight=240mm,left=20mm,right=20mm, top=25mm, bottom=15mm]{geometry}               %定义版面paperwidth=185mm,paperheight=230mm
\usepackage[CJKbookmarks, colorlinks, bookmarksnumbered=true,pdfstartview=FitH,linkcolor=black]{hyperref}   % 选项去掉链接红色方框, 去掉书签可得到想要的印刷效果
%\usepackage{etex}                           %定义更多的计数器,否则警告"no room for a new \count"
%-------------------------------------------------------------------------------------------------------
%\usepackage{graphicx}
%\usepackage{picins}                         %图文混排(文包图),可带方框
%\usepackage{picinpar}                  %图文混排(文包图)  picins在前,picinpar在后,否则有问题
%\usepackage{overpic}                    %加背景图等
%-------------------------------------------------------------------------
\usepackage{fancyhdr}                    % 页眉和页脚的相关定义
\usepackage[perpage]{footmisc}    % 脚注控制
%\usepackage[perpage,symbol]{footmisc}    % 脚注控制
\usepackage{titlesec}                        % 控制标题的宏包
\usepackage{titletoc}                         % 控制目录的宏包
\usepackage{caption2}                       %浮动图形和表格标题样式
\usepackage{setspace}                       % 定制表格和图形的多行标题行距，行间距的宏包
%------------------------------------------------------------------------
\usepackage[usenames,dvipsnames]{color}   % 支持彩色
\usepackage{xcolor}               %单元格加背景
\usepackage{multirow}           %制作表格主要用到\multicolumn，\multirow和\cline,其中，要使用\multirow
%--------------------------------------------------------------------------------------------------------------------------
%\usepackage{everb}                  % 扩展了verbatim，可以加边框、前景色，背景色和行号等，\begin{colorboxed} 这里的内容可以跨页\end{colorboxed}
%\usepackage{listings}          % 可对关键词、注释和字符串等使用不同的字体和颜色，也可以为代码添加边框、背景等风格
\usepackage{fancybox}           % 边框,有阴影,fancybox提供了五种式样\fbox,\shadowbox,\doublebox,\ovalbox,\Ovalbox。
%\usepackage[listings,theorems]{tcolorbox}                % 用于定义类似于beamer的环境框
\usepackage{float}
\usepackage{multicol}           %用于实现在同一页中实现不同的分栏
%\usepackage{flushend}       %分两栏后,\flushend 命令使双栏底部对齐,\raggedend 命令则取消底部对齐,文档最后有效
%\usepackage{balance}           %分两栏后,\balance  命令使双栏底部对齐,\balance   命令则取消底部对齐,任何位置有效

%--------------------------------------------------------------------------------------------
\usepackage{extarrows}          %实现长等号等  This package provides some extensible arrows in maths. mode    (more than \xlongequal \xleftarrow \xrightarrow \xLongleftarrow \xLongrightarrow \xLongleftrightarrow
% wasysym 会在小字号下产生 U/wasy 字形替换警告；本书正文没有使用其特殊符号。
%---------------------------------------------------------------------------------
\usepackage{indentfirst}               % 首行缩进宏包
%\usepackage{color}                        % 支持彩色
%------------------------------------------------------------------------------------
% XeLaTeX 下不要使用旧的 times/helvet 字体包，否则会产生 TU/ptm 与 TU/phv 字形警告。
\IfFontExistsTF{TeX Gyre Termes}
  {\setmainfont{TeX Gyre Termes}
   \setsansfont{TeX Gyre Heros}
   \setmonofont{TeX Gyre Cursor}}
  {}
\usepackage{amssymb}            % AMSLaTeX宏包 用来排出更加漂亮的公式
\usepackage{mathrsfs}               % 不同于\mathcal or \mathfrak 之类的英文花体字体
\usepackage{bm}                         % 处理数学公式中的黑斜体的宏包

\usepackage{tikz,pifont,diagbox}   % 带圆圈的数字，例如\circled{4} 
\newcommand{\circled}[1]{%
  \tikz[baseline=(char.base)]{
    \node[shape=circle,draw,inner sep=1pt,line width=0.5pt] (char) {\small #1};
  }%
}

% 清理日志：旧稿中少量中文标点/文字会落入 Latin 字体或数学模式。
\newunicodechar{，}{\ifmmode\text{\CJKfamily{zhsong}，}\else{\CJKfamily{zhsong}，}\fi}
\newunicodechar{尼}{\ifmmode\text{\CJKfamily{zhsong}尼}\else{\CJKfamily{zhsong}尼}\fi}
\newunicodechar{姆}{\ifmmode\text{\CJKfamily{zhsong}姆}\else{\CJKfamily{zhsong}姆}\fi}
\newunicodechar{≤}{\ensuremath{\leq}}
\hbadness=10000
\vbadness=10000
\hfuzz=80pt
\vfuzz=80pt

\usepackage[amsmath,thmmarks]{ntheorem}         % 定理类环境宏包，其中 amsmath 选项用来兼容 AMS LaTeX 的宏包thmmarks 选项，可以自动恰当地放置定理类环境的结束标记；它还能像图形目录那样生成定理类环境目录
\usepackage{graphicx}
\usepackage{epsfig}
\usepackage{booktabs}
\usepackage{rotating}%将表格或者图片处理成单独一个页面
\usepackage{imakeidx}%可生成索引的标准 LaTeX 宏包
\usepackage{pdfpages}%在LaTeX里插入全页的pdf

%----------------------------------------------------------------------------------------------------------
%\usepackage{flafter}                               % 因为图形可浮动到当前页的顶部， 要防止这种情况，可使用 flafter 宏包
%\usepackage[below]{placeins}
%\usepackage{floatflt}                              % 图文混排用宏包
%\usepackage{rotating}                            % 图形和表格的控制
%\usepackage{type1cm}                           % tex1cm宏包，控制字体的大小
%\usepackage{endfloat}                            %可将浮动对象放置到文件的最后
%----------------------------------------------------------------------------------------------------------------------
\usepackage{fancyref}                                % 支持引用的宏包
%\usepackage{cite}                                    % 支持引用缩写的宏包
%\usepackage{natbib}
\makeatletter
\@ifundefined{CJKnumber}
  {\def\CJKnumber#1{\ifcase#1\or
                    1\or2\or3\or4\or5\or
                    6\or7\or8\or9\or10\or11\or12\or13\or14\or15\or16\or17\or18\or19\or20\fi}}{}
\pdfstringdefDisableCommands{%
  \def\CJKnumber#1{#1}%
  \def\CJKprechaptername{第}%
  \def\thechapter{}%
  \def\CJKthechapter{}%
  \def\CJKchaptername{章}%
  \def\color#1{}%
  \def\hei{}%
  \def\song{}%
}

%\usepackage{amsthm}
 

\renewcommand\contentsname{目\qquad 录}
\renewcommand\listfigurename{插图}
\renewcommand\listtablename{表格}

\@ifundefined{chapter}
  {\renewcommand\refname{参考文献}}
  {\renewcommand\bibname{参考文献}}

\renewcommand\indexname{索引}

\renewcommand\figurename{图}

\renewcommand\tablename{表}

\newcommand\CJKprepartname{第}
\newcommand\CJKpartname{部分}
\newcommand\CJKthepart{\CJKnumber{\@arabic\c@part}}

\@ifundefined{chapter}{}{
  \newcommand\CJKprechaptername{第}
  \newcommand\CJKchaptername{章}
  \newcommand\CJKthechapter{\CJKnumber{\@arabic\c@chapter}}}
\AtBeginDocument{%
  \pdfstringdefDisableCommands{%
    \def\CJKnumber#1{#1}%
    \def\CJKprechaptername{}%
    \def\CJKthechapter{}%
    \def\CJKchaptername{}%
  }%
}
  

\renewcommand\appendixname{附录~\@Alph\c@chapter}

\@ifundefined{mainmatter}
  {\renewcommand\abstractname{摘要}}{}

% \renewcommand\ccname{}                     %   ?
% \renewcommand\enclname{附件}
% \newcommand\prepagename{}                  %   ?
% \newcommand\postpagename{}                 %   ?
% \renewcommand\headtoname{}                 %   ?
% \renewcommand\seename{}                    %   ?

\let\CJK@todaysave=\today
\def\CJK@todaysmall{~\the\year~年~\the\month~月~\the\day~日}
\def\CJK@todaybig{\CJKdigits{\the\year}年\CJKnumber{\the\month}月\CJKnumber{\the\day}日}
\def\CJK@today{\CJK@todaysmall}
\renewcommand\today{\CJK@today}
\newcommand\CJKtoday[1][1]{%
  \ifcase#1\def\CJK@today{\CJK@todaysave}
  \or\def\CJK@today{\CJK@todaysmall}
  \or\def\CJK@today{\CJK@todaybig}
  \fi}

%
% modify the definitions of Standard LaTeX document class
%
\@ifundefined{chapter}{
  \def\@part[#1]#2{%
      \ifnum \c@secnumdepth >\m@ne
        \refstepcounter{part}%
%       \addcontentsline{toc}{part}{\thepart\hspace{1em}#1}%
        \addcontentsline{toc}{part}{\CJKprepartname\expandafter\noexpand\CJKthepart\CJKpartname\hspace{1em}#1}%
      \else
        \addcontentsline{toc}{part}{#1}%
      \fi
      {\parindent \z@ \raggedright
       \interlinepenalty \@M
       \normalfont
       \ifnum \c@secnumdepth >\m@ne
%        \Large\bfseries \partname\nobreakspace\thepart
         \Large\bfseries \CJKprepartname\CJKthepart\CJKpartname
         \par\nobreak
       \fi
       \huge \bfseries #2%
       \markboth{}{}\par}%
      \nobreak
      \vskip 3ex
      \@afterheading}
}{
  \def\@part[#1]#2{%
      \ifnum \c@secnumdepth >-2\relax
        \refstepcounter{part}%
%       \addcontentsline{toc}{part}{\thepart\hspace{1em}#1}%
        \addcontentsline{toc}{part}{\CJKprepartname\expandafter\noexpand\CJKthepart\CJKpartname\hspace{1em}#1}%
      \else
        \addcontentsline{toc}{part}{#1}%
      \fi
      \markboth{}{}%
      {\centering
       \interlinepenalty \@M
       \normalfont
       \ifnum \c@secnumdepth >-2\relax
%        \huge\bfseries \partname\nobreakspace\thepart
         \huge\bfseries \CJKprepartname\CJKthepart\CJKpartname
         \par
         \vskip 20\p@
       \fi
       \Huge \bfseries #2\par}%
      \@endpart}
  \if@twoside
    \def\chaptermark#1{%
      \markboth {\MakeUppercase{%
        \ifnum \c@secnumdepth >\m@ne
          \if@mainmatter
%           \@chapapp\ \thechapter. \ %
            \CJKprechaptername\CJKthechapter\CJKchaptername \ %
          \fi
        \fi
        #1}}{}}%
    \def\sectionmark#1{%
      \markright {\MakeUppercase{%
        \ifnum \c@secnumdepth >\z@
%         \thesection. \ %
          \thesection \ %
        \fi
        #1}}}
  \else
    \def\chaptermark#1{%
      \markright {\MakeUppercase{%
        \ifnum \c@secnumdepth >\m@ne
          \if@mainmatter
%           \@chapapp\ \thechapter. \ %
            \CJKprechaptername\CJKthechapter\CJKchaptername \ %
          \fi
        \fi
        #1}}}
  \fi
  \def\@chapter[#1]#2{\ifnum \c@secnumdepth >\m@ne
                         \if@mainmatter
                           \refstepcounter{chapter}%
%                          \typeout{\@chapapp\space\thechapter.}%
                           \typeout{\CJKprechaptername\CJKthechapter\CJKchaptername}%
                           \addcontentsline{toc}{chapter}%
%                                    {\protect\numberline{\thechapter}#1}%
                                     {\texorpdfstring{\protect\numberline{}\CJKprechaptername%
                                      \expandafter\noexpand\CJKthechapter\CJKchaptername\hspace{0.8em}#1}{#1}}%
                         \else
                           \addcontentsline{toc}{chapter}{#1}%
                         \fi
                      \else
                        \addcontentsline{toc}{chapter}{#1}%
                      \fi
                      \chaptermark{#1}%
                      \addtocontents{lof}{\protect\addvspace{10\p@}}%
                      \addtocontents{lot}{\protect\addvspace{10\p@}}%
                      \if@twocolumn
                        \@topnewpage[\@makechapterhead{#2}]%
                      \else
                        \@makechapterhead{#2}%
                        \@afterheading
                      \fi}
  \def\@makechapterhead#1{%
    \vspace*{50\p@}%
    {\parindent \z@ \raggedright \normalfont
      \ifnum \c@secnumdepth >\m@ne
        \if@mainmatter
%         \huge\bfseries \@chapapp\space \thechapter
          \huge\bfseries \CJKprechaptername\CJKthechapter\CJKchaptername
          \par\nobreak
          \vskip 20\p@
        \fi
      \fi
      \interlinepenalty\@M
      \Huge \bfseries #1\par\nobreak
      \vskip 40\p@
    }}
  \renewcommand*\l@chapter[2]{%
    \ifnum \c@tocdepth >\m@ne
      \addpenalty{-\@highpenalty}%
      \vskip 1.0em \@plus\p@
%     \setlength\@tempdima{1.5em}%
      \setlength\@tempdima{0em}%
      \begingroup
        \parindent \z@ \rightskip \@pnumwidth
        \parfillskip -\@pnumwidth
        \leavevmode \bfseries
        \advance\leftskip\@tempdima
        \hskip -\leftskip
        #1\nobreak\hfil \nobreak\hb@xt@\@pnumwidth{\hss #2}\par
        \penalty\@highpenalty
      \endgroup
    \fi}
 

\let\@appendix\appendix
\renewcommand\appendix{\@appendix%
  \def\CJKprechaptername{\relax}%
  \def\CJKthechapter{\relax}%
  \def\CJKchaptername{\appendixname}}
}  %end of \@ifundefined{chapter}

\def\numberline#1{\ifdim\@tempdima>0pt%
  \settowidth\@tempdimb{#1\space}%
  \ifdim\@tempdima<\@tempdimb%
    \@tempdima=\@tempdimb%
  \fi%
  \hb@xt@\@tempdima{#1\hfil}%
  \fi}
\long\def\@makefntext#1{\leavevmode\hbox{}\kern1em $^{\@thefnmark}$\footnotesize#1}
\fboxsep = 3pt
\fboxrule = .4pt
% Speciali footnotetext pirmam lapui
\def\sfootnotetext[#1]{\begingroup\let\protect\noexpand
      \xdef\@thefnmark{#1}\endgroup
    \@footnotetext}
\makeatother 


\usepackage{natbib}
\bibliographystyle{apalike}

\hypersetup{
  colorlinks=true,
  citecolor=blue,
  linkcolor=blue,
  urlcolor=cyan,
  hypertexnames=false,
  pdftitle={Overleaf Example},
}

\usepackage{amsmath}
\usepackage[most]{tcolorbox}
\newtcolorbox{mybox}[1]{
   colback=white,
    colframe=black,
    colbacktitle=gray!30,
    coltitle=black,
    fonttitle=\bfseries,
    title=#1,
    % --- 关键排版设置 ---
    float=!ht,  % ! 表示强制尝试，“h”当前位置，“t”页顶
}

\newtheorem{exercise}{习题}[chapter]
